\subsection{Held common-appearance results}
\label{sec:coupled-appearance-results}
Table~\ref{tab:coupled-appearance} reports all 48 queries for each fully covered method. All methods transport the same Gaussian appearance using the same rendering pipeline; this isolates image differences induced by predicted motion and is not a comparison of native end-to-end reconstruction systems.
\begin{table}[t]
\centering\small
\caption{Six-family held common-appearance evaluation. Foreground-union PSNR and SSIM and image LPIPS use noninitial frames. Means average the same 48 queries per row. Higher PSNR/SSIM and lower LPIPS are better.}
\label{tab:coupled-appearance}
\begin{tabular}{lrrr}
\toprule Method & PSNR & SSIM & LPIPS \\
\midrule
Ours & 35.61 & 0.877 & 0.0101 \\
Analytic control & 36.02 & 0.887 & 0.0093 \\
PhysTwin adaptation & 25.68 & 0.731 & 0.0404 \\
Spring-Gaus adaptation & 25.88 & 0.721 & 0.0425 \\
\bottomrule
\end{tabular}
\end{table}
Across 24 independent worlds, the paired foreground PSNR difference (ours minus baseline) is $+9.93$\,dB against PhysTwin (world-bootstrap 95\% interval $[8.24,11.76]$) and $+9.73$\,dB against Spring-Gaus ($[7.93,11.67]$). The analytic control is slightly better: our difference is $-0.41$\,dB ($[-0.77,-0.08]$). LPIPS differences favor ours over the two native adaptations, while favoring the analytic control over ours. These comparisons are specific to the declared adapters, observation prefixes and fitting budgets.

UniPhy completes one of eight supported drop queries; seven fail with a native CFL error. The other 40 slots are unsupported by this adapter. Because no UniPhy world completes both queries, no paired-world image comparison is reported. Its one completed query has foreground PSNR $12.24$\,dB; this isolated descriptive value is not comparable to the 48-query averages above.

Physical-error evaluation and the two remaining spatial-reference refinement runs are still in progress at this draft stage. Image fidelity alone does not establish mechanical accuracy.

\begin{table}[t]
\centering\scriptsize
\caption{Foreground PSNR by scene family (dB; eight queries each). Averages describe the frozen held cohort, not population guarantees.}
\label{tab:coupled-appearance-family}
\begin{tabular}{lrrrr}
\toprule Family & Ours & Analytic & Spring & PhysTwin \\
\midrule
Clamps/release & 28.78 & 28.96 & 19.91 & 20.06 \\
Curved pushing & 30.63 & 30.20 & 24.21 & 21.86 \\
Container drop & 45.37 & 45.83 & 27.78 & 27.82 \\
Rigid impact & 40.70 & 42.73 & 36.94 & 35.72 \\
Two soft bodies & 40.93 & 41.16 & 27.16 & 27.50 \\
Heterogeneous & 27.26 & 27.23 & 19.30 & 21.14 \\
\bottomrule
\end{tabular}
\end{table}

\subsection{Held comparison media}
The six-family media set contains twelve clips: the first predefined world
(world00) in every family, with both held interactions. This selection does
not depend on prediction errors. Each clip shows six panels: the simulation
reference, ours, the analytic control, Spring-Gaus, PhysTwin and UniPhy.
Unsupported and failed slots are labeled explicitly. Complete predictions
use synchronized saved physical times, the same orbiting camera and common
axis limits; deformation is not amplified. Playback is 24 frames per second
for states sampled every 2\,ms, corresponding to approximately $20.8\times$
slow motion. These checkerboard surface visualizations expose geometry and
interaction differences; they are separate from the common-Gaussian image
evaluation in Tables~\ref{tab:coupled-appearance}
and~\ref{tab:coupled-appearance-family}.
The local media filenames follow
\texttt{held-\{family\}-world00-query\{0,1\}.mp4}.
The accompanying audit records source trajectory hashes, movie hashes,
frame counts and successful full-stream decoding for all twelve clips.
